\section{Desargues: the identity and both implications}\label{sec:des}

\begin{lemma}[The universal Desargues identity]\label{lem:des-id}
With the notation of Theorem~\ref{thm:des-summary}, over every
commutative ring,
\begin{equation}\label{eq:des-id}
 \boxed{\br X Y Z=\br a b c\,\br A B C\,
              \br{a\times A}{b\times B}{c\times C}.}
\end{equation}
\end{lemma}
\begin{proof}
By \eqref{eq:doublecross},
\begin{align*}
 X&=\br a b B\,A-\br a b A\,B,\\
 Y&=\br b c C\,B-\br b c B\,C,\\
 Z&=\br c a A\,C-\br c a C\,A.
\end{align*}
Multilinearity gives eight terms. Six have repeated columns, leaving
\[
 \br X Y Z=E\br A B C,
\]
\[
 E=\br a b B\br b c C\br c a A
       -\br a b A\br b c B\br c a C.
\]
Let $T=(a\ b\ c)$ and $J=(a\times A\ \ b\times B\ \ c\times C)$.
Their matrix product is
\[
 T^{\mathsf t}J=
 \begin{pmatrix}
 0&\br a b B&-\br c a C\\
 -\br a b A&0&\br b c C\\
 \br c a A&-\br b c B&0
 \end{pmatrix}.
\]
Its diagonal is zero; the two surviving permutation terms in its
determinant are exactly those defining $E$. Multiplicativity gives
$E=\det T\det J$, proving \eqref{eq:des-id}. No cancellation or
nondegeneracy assumption entered this identity.
\end{proof}

\begin{proof}[Proof of Theorem~\ref{thm:des-summary}]
Suppose a nonzero center $o$ is given. Its three incidences and
Lemma~\ref{lem:inc-det} give $\br{j_1}{j_2}{j_3}=0$.
Equation~\eqref{eq:des-id} gives $\br X Y Z=0$.
Since $X\ne0$, Lemma~\ref{lem:axis} constructs the nonzero
axis $h=\axis(X,Y,Z)$. This proves the forward implication.

Conversely suppose a nonzero line $h$ through $X,Y,Z$ is given.
Lemma~\ref{lem:inc-det}, with $h$ in place of $o$, gives
$\br X Y Z=0$. Both $\br a b c$ and $\br A B C$ are nonzero,
so two applications of integer cancellation in \eqref{eq:des-id}
give $\br{j_1}{j_2}{j_3}=0$. Since $j_1\ne0$, the same algorithm
applied dually produces $o=\axis(j_1,j_2,j_3)$, a nonzero common center.
Every existential introduction uses this displayed term; every
existential elimination merely opens the supplied witness.
\end{proof}

\subsection{An exact example}
Take
\[
 a=(1,0,1),\quad b=(0,1,1),\quad c=(-1,-1,1),
\]
\[
 A=(1,0,2),\quad B=(0,1,3),\quad C=(-1,-1,4).
\]
The triangle determinants are $3$ and $9$. The joining lines are
$(0,-1,0)$, $(2,0,0)$, and $(-3,3,0)$, all through $(0,0,1)$.
The side intersections computed by the cross products are
\[
 X=(2,-1,1),\quad Y=(2,5,1),\quad Z=(-4,-1,-2).
\]
All lie on $(1,0,-2)$, or $x=2$ in the chart $z=1$.
The algorithm returns $X\times Y=(-6,0,12)$, representing that line.

\subsection{Why the converse needs nondegenerate triangles}
For
\[
 a=(1,0,0),\ b=(0,1,0),\ c=(1,1,0),\qquad
 A=(0,0,1),\ B=(1,0,1),\ C=(0,1,1),
\]
the side intersections are $(-1,0,0),(-1,1,0),(0,-1,0)$,
so are collinear. All joins and meets are nonzero. But the
joining-line determinant is $2$, so Lemma~\ref{lem:inc-det}
rules out a common nonzero center. Here $\br a b c=0$.
This identifies precisely why the converse cannot discard the
hypotheses that allow the triangle determinants to be canceled.

